\documentclass[12pt,a4paper]{article}

% -------------------------------------------------------
% Pacotes
% -------------------------------------------------------
\usepackage[utf8]{inputenc}
\usepackage[T1]{fontenc}
\usepackage[portuguese]{babel}
\usepackage{amsmath,amssymb,amsthm}
\usepackage{geometry}
\usepackage{hyperref}
\usepackage{graphicx}
\usepackage{enumitem}
\usepackage{booktabs}
\usepackage{xcolor}
\usepackage{fancyhdr}
\usepackage{titlesec}
\usepackage{mdframed}
\usepackage{listings}
\usepackage{mathtools}

\geometry{margin=2.5cm}

% -------------------------------------------------------
% Ambientes de teoremas
% -------------------------------------------------------
\newtheorem{theorem}{Teorema}[section]
\newtheorem{lemma}[theorem]{Lema}
\newtheorem{corollary}[theorem]{Corolário}
\newtheorem{proposition}[theorem]{Proposição}
\theoremstyle{definition}
\newtheorem{definition}[theorem]{Definição}
\newtheorem{example}[theorem]{Exemplo}
\newtheorem{exercise}{Exercício}[section]
\theoremstyle{remark}
\newtheorem{remark}[theorem]{Observação}

% -------------------------------------------------------
% Caixas coloridas para resultados principais
% -------------------------------------------------------
\newmdenv[linecolor=blue!60,backgroundcolor=blue!5,
          roundcorner=4pt,innerleftmargin=10pt,innerrightmargin=10pt]{keybox}

\newmdenv[linecolor=green!50!black,backgroundcolor=green!5,
          roundcorner=4pt,innerleftmargin=10pt,innerrightmargin=10pt]{exbox}

% -------------------------------------------------------
\begin{document}

\tableofcontents
\newpage

% ================================================================
%  CAPÍTULO 2 -- CÁLCULO INTEGRAL EM R
% ================================================================
\section{Cálculo Integral em \texorpdfstring{$\mathbb{R}$}{R}}

Este capítulo segue os tópicos do programa 2.1--2.5 especificados para os
cursos de EI/EE e EM. Abrange primitivas e as suas técnicas de cálculo,
o integral definido, o Teorema Fundamental do Cálculo Integral e
aplicações ao cálculo de áreas, volumes e comprimentos de arco.

% -----------------------------------
\subsection{Primitivas e Técnicas de Primitivação (2.1)}
% -----------------------------------

\begin{definition}[Primitiva]
  Uma função $F : I \to \mathbb{R}$ diz-se uma \emph{primitiva} (ou
  \emph{antiderivada}) de $f : I \to \mathbb{R}$ num intervalo aberto $I$ se
  \[
    F'(x) = f(x) \quad \forall\, x \in I.
  \]
  A família de todas as primitivas escreve-se $\displaystyle\int f(x)\,dx = F(x)+C$,
  onde $C \in \mathbb{R}$ é uma constante arbitrária.
\end{definition}

\begin{remark}
  Duas primitivas da mesma função diferem por uma constante, pelo que o
  integral indefinido é único a menos de $C$.
\end{remark}

\subsubsection*{Tabela de Primitivas Imediatas}

\begin{center}
\begin{tabular}{cc}
  \toprule
  $f(x)$ & $\displaystyle\int f(x)\,dx$ \\
  \midrule
  $x^n\ (n\neq-1)$ & $\dfrac{x^{n+1}}{n+1}+C$ \\[6pt]
  $\dfrac{1}{x}$ & $\ln|x|+C$ \\[6pt]
  $e^x$ & $e^x+C$ \\[6pt]
  $\sin x$ & $-\cos x+C$ \\[6pt]
  $\cos x$ & $\sin x+C$ \\[6pt]
  $\sec^2 x$ & $\tan x+C$ \\[6pt]
  $\dfrac{1}{1+x^2}$ & $\arctan x + C$ \\[6pt]
  $\dfrac{1}{\sqrt{1-x^2}}$ & $\arcsin x + C$ \\
  \bottomrule
\end{tabular}
\end{center}

\subsubsection*{Primitivação por Partes}

\begin{keybox}
\textbf{Fórmula:}
\[
  \int u\,dv = uv - \int v\,du.
\]
Escolhe-se $u$ e $dv$ de modo que $\int v\,du$ seja mais simples. O
mnemónico \textbf{LIATE} (Logarítmica, Inversa trigonométrica, Algébrica,
Trigonométrica, Exponencial) sugere a prioridade na escolha de $u$.
\end{keybox}

\begin{example}[Primitivação por Partes]
  Calcule $\displaystyle\int x e^x\,dx$.

  \textbf{Resolução.} Tome-se $u=x$, $dv=e^x\,dx$. Então $du=dx$, $v=e^x$.
  \[
    \int x e^x\,dx = xe^x - \int e^x\,dx = xe^x - e^x + C = e^x(x-1)+C.
  \]
\end{example}

\subsubsection*{Primitivação por Substituição}

\begin{keybox}
\textbf{Fórmula:} Se $x = g(t)$, então
\[
  \int f(x)\,dx = \int f(g(t))\,g'(t)\,dt.
\]
\end{keybox}

\begin{example}[Substituição]
  Calcule $\displaystyle\int \frac{2x}{x^2+1}\,dx$.

  \textbf{Resolução.} Seja $u = x^2+1$, logo $du = 2x\,dx$.
  \[
    \int \frac{2x}{x^2+1}\,dx = \int \frac{du}{u} = \ln|u|+C = \ln(x^2+1)+C.
  \]
\end{example}

\subsubsection*{Potências de Funções Trigonométricas}

Para integrar $\sin^m x\cos^n x$:
\begin{itemize}[nosep]
  \item Se $m$ for ímpar: substitua $u=\cos x$.
  \item Se $n$ for ímpar: substitua $u=\sin x$.
  \item Se ambos forem pares: use as identidades do ângulo duplo
        $\sin^2 x = \tfrac{1-\cos 2x}{2}$,
        $\cos^2 x = \tfrac{1+\cos 2x}{2}$.
\end{itemize}

\begin{example}
  Calcule $\displaystyle\int \sin^3 x\,dx$.

  \textbf{Resolução.} Escreva $\sin^3 x = \sin x(1-\cos^2 x)$; seja $u=\cos x$:
  \[
    \int \sin^3 x\,dx = -\int (1-u^2)\,du = -u + \tfrac{u^3}{3}+C
      = -\cos x + \tfrac{\cos^3 x}{3}+C.
  \]
\end{example}

\subsubsection*{Frações Racionais (Decomposição em Frações Parciais)}

Toda a função racional própria $\frac{P(x)}{Q(x)}$ pode ser decomposta em
frações parciais. Para uma raíz real simples $r$ de $Q$:
\[
  \frac{A}{x-r}.
\]
Para um fator quadrático irredutível $x^2+px+q$:
\[
  \frac{Bx+C}{x^2+px+q}.
\]

\begin{example}
  Calcule $\displaystyle\int \frac{1}{x^2-1}\,dx$.

  \textbf{Resolução.} Decomposição: $\frac{1}{x^2-1}=\frac{1}{2}\!\left(\frac{1}{x-1}-\frac{1}{x+1}\right)$.
  \[
    \int \frac{dx}{x^2-1} = \frac{1}{2}\ln\left|\frac{x-1}{x+1}\right|+C.
  \]
\end{example}

% -----------------------------------
\subsection{O Integral Definido (2.2)}
% -----------------------------------

\begin{definition}[Integral de Riemann]
  Seja $f:[a,b]\to\mathbb{R}$ limitada. Para uma partição
  $P=\{a=x_0<x_1<\cdots<x_n=b\}$, a \emph{soma de Riemann} é
  \[
    S(f,P) = \sum_{i=1}^n f(c_i)(x_i - x_{i-1}),\quad c_i\in[x_{i-1},x_i].
  \]
  Se $S(f,P)$ convergir para um limite $L$ quando a norma da partição
  $\|P\|\to 0$ (independentemente da escolha de $c_i$), então $f$ é
  \emph{integrável no sentido de Riemann} e $\displaystyle\int_a^b f(x)\,dx = L$.
\end{definition}

\begin{keybox}
\textbf{Propriedades principais:}
\begin{align*}
  &\int_a^b [\alpha f+\beta g]\,dx = \alpha\int_a^b f\,dx + \beta\int_a^b g\,dx
    \quad (\text{linearidade})\\
  &\int_a^b f\,dx = \int_a^c f\,dx + \int_c^b f\,dx
    \quad (\text{aditividade em } [a,c]\cup[c,b])\\
  &f(x)\geq 0 \implies \int_a^b f\,dx \geq 0
    \quad (\text{positividade})
\end{align*}
\end{keybox}

% -----------------------------------
\subsection{O Teorema Fundamental do Cálculo Integral (2.3)}
% -----------------------------------

\begin{theorem}[Teorema Fundamental do Cálculo Integral]
  Seja $f:[a,b]\to\mathbb{R}$ contínua.
  \begin{enumerate}
    \item \textbf{(Parte 1)} A função $G(x)=\displaystyle\int_a^x f(t)\,dt$
          é diferenciável em $(a,b)$ e $G'(x)=f(x)$.
    \item \textbf{(Parte 2)} Se $F$ for qualquer primitiva de $f$, então
          \[
            \int_a^b f(x)\,dx = F(b)-F(a) =: \Big[F(x)\Big]_a^b.
          \]
  \end{enumerate}
\end{theorem}

\begin{example}
  Calcule $\displaystyle\int_0^{\pi} \sin x\,dx$.

  \textbf{Resolução.} Uma primitiva de $\sin x$ é $-\cos x$.
  \[
    \int_0^{\pi}\sin x\,dx = \Big[-\cos x\Big]_0^{\pi}
      = (-\cos\pi)-(-\cos 0) = 1+1 = 2.
  \]
\end{example}

% -----------------------------------
\subsection{Aplicações e Comprimento do Arco (2.4 e 2.5)}
% -----------------------------------

\subsubsection*{Área entre curvas}

Se $f(x)\geq g(x)$ em $[a,b]$, a área da região delimitada é:
\[
  A = \int_a^b \bigl[f(x)-g(x)\bigr]\,dx.
\]

\begin{example}
  Determine a área entre $f(x)=x^2$ e $g(x)=x$ em $[0,1]$.

  \textbf{Resolução.} Como $x\geq x^2$ em $[0,1]$:
  \[
    A = \int_0^1(x-x^2)\,dx = \left[\frac{x^2}{2}-\frac{x^3}{3}\right]_0^1
      = \frac{1}{2}-\frac{1}{3} = \frac{1}{6}.
  \]
\end{example}

\subsubsection*{Volume de um Sólido de Revolução (Método dos Discos)}

Ao rodar $y=f(x)\geq 0$ em torno do eixo $Ox$ de $a$ a $b$:
\[
  V = \pi\int_a^b [f(x)]^2\,dx.
\]

\begin{example}
  Volume da esfera de raio $r$: roda-se $f(x)=\sqrt{r^2-x^2}$ em
  $[-r,r]$ em torno do eixo $Ox$.

  \textbf{Resolução.}

  \textbf{Passo 1 -- Identificar a função.} A semicircunferência superior de
  raio $r$ é dada por $f(x)=\sqrt{r^2-x^2}$, com $x\in[-r,r]$. Ao rodar
  este arco em torno do eixo $Ox$ obtém-se uma esfera completa.

  \textbf{Passo 2 -- Aplicar a fórmula dos discos.}
  \[
    V = \pi\int_{-r}^r [f(x)]^2\,dx = \pi\int_{-r}^r (r^2-x^2)\,dx.
  \]

  \textbf{Passo 3 -- Calcular o integral.}
  \[
    V = \pi\left[r^2 x - \frac{x^3}{3}\right]_{-r}^r
      = \pi\left[\left(r^3-\frac{r^3}{3}\right)-\left(-r^3+\frac{r^3}{3}\right)\right]
      = \pi\cdot 2\left(r^3-\frac{r^3}{3}\right)
      = \pi\cdot\frac{4r^3}{3} = \frac{4}{3}\pi r^3.
  \]
\end{example}

\subsubsection*{Comprimento do Arco (2.5)}

O comprimento da curva $y=f(x)$ em $[a,b]$:
\[
  L = \int_a^b \sqrt{1+[f'(x)]^2}\,dx.
\]

Para uma curva paramétrica $(x(t),y(t))$, $t\in[\alpha,\beta]$:
\[
  L = \int_\alpha^\beta \sqrt{[\dot x(t)]^2+[\dot y(t)]^2}\,dt.
\]

\begin{example}
  Determine o comprimento de $y=\frac{2}{3}x^{3/2}$ de $x=0$ a $x=3$.

  \textbf{Resolução.} $f'(x)=x^{1/2}$, logo $1+[f']^2=1+x$.
  \[
    L = \int_0^3\sqrt{1+x}\,dx = \left[\tfrac{2}{3}(1+x)^{3/2}\right]_0^3
      = \tfrac{2}{3}(8-1) = \tfrac{14}{3}.
  \]
\end{example}

% -----------------------------------
\subsection{Exercícios de Prática -- Capítulo 2}
% -----------------------------------

\begin{exercise}
  Calcule os seguintes integrais indefinidos:
  \begin{enumerate}[label=(\alph*)]
    \item $\displaystyle\int (3x^2-2x+5)\,dx$
    \item $\displaystyle\int \frac{x}{x^2+4}\,dx$
    \item $\displaystyle\int x^2 e^x\,dx$
    \item $\displaystyle\int \frac{x+1}{x^2+x-2}\,dx$
    \item $\displaystyle\int \cos^2 x\,dx$
  \end{enumerate}

  \medskip\textbf{Resolução.}

  \textbf{(a)} Aplica-se a linearidade e a regra da potência:
  \[
    \int (3x^2-2x+5)\,dx = 3\cdot\frac{x^3}{3} - 2\cdot\frac{x^2}{2} + 5x + C
    = x^3 - x^2 + 5x + C.
  \]

  \textbf{(b)} Reconhece-se a forma $\int \frac{f'(x)}{f(x)}\,dx$. Seja $u=x^2+4$,
  $du=2x\,dx$, logo $x\,dx=\frac{du}{2}$:
  \[
    \int \frac{x}{x^2+4}\,dx = \frac{1}{2}\int\frac{du}{u}
    = \frac{1}{2}\ln|u|+C = \frac{1}{2}\ln(x^2+4)+C.
  \]

  \textbf{(c)} Aplica-se primitivação por partes duas vezes (LIATE: $u=x^2$,
  $dv=e^x\,dx$):
  \begin{align*}
    &\textbf{1.ª vez:}\quad u=x^2,\ dv=e^x\,dx \Rightarrow du=2x\,dx,\ v=e^x.\\
    &\int x^2 e^x\,dx = x^2 e^x - 2\int x e^x\,dx.\\
    &\textbf{2.ª vez:}\quad u=x,\ dv=e^x\,dx \Rightarrow du=dx,\ v=e^x.\\
    &\int x e^x\,dx = xe^x - e^x + C.\\
    &\textbf{Resultado:}\quad \int x^2 e^x\,dx = x^2 e^x - 2(xe^x - e^x)+C
      = e^x(x^2-2x+2)+C.
  \end{align*}

  \textbf{(d)} Fatoriza-se o denominador: $x^2+x-2=(x+2)(x-1)$.
  Decomposição em frações parciais:
  \[
    \frac{x+1}{(x+2)(x-1)} = \frac{A}{x+2}+\frac{B}{x-1}.
  \]
  Multiplicando ambos os membros por $(x+2)(x-1)$: $x+1 = A(x-1)+B(x+2)$.
  \begin{itemize}[nosep]
    \item $x=1$: $2=3B \Rightarrow B=\tfrac{2}{3}$.
    \item $x=-2$: $-1=-3A \Rightarrow A=\tfrac{1}{3}$.
  \end{itemize}
  \[
    \int \frac{x+1}{x^2+x-2}\,dx = \frac{1}{3}\ln|x+2|+\frac{2}{3}\ln|x-1|+C.
  \]

  \textbf{(e)} Usa-se a identidade do ângulo duplo: $\cos^2 x = \dfrac{1+\cos 2x}{2}$.
  \[
    \int \cos^2 x\,dx = \int \frac{1+\cos 2x}{2}\,dx
    = \frac{x}{2} + \frac{\sin 2x}{4} + C.
  \]
\end{exercise}

\begin{exercise}
  Calcule os seguintes integrais definidos:
  \begin{enumerate}[label=(\alph*)]
    \item $\displaystyle\int_1^e \ln x\,dx$
    \item $\displaystyle\int_0^1 x\sqrt{1-x^2}\,dx$
    \item $\displaystyle\int_0^{\pi/2}\sin^2 x\,dx$
  \end{enumerate}

  \medskip\textbf{Resolução.}

  \textbf{(a)} Primitivação por partes: $u=\ln x$, $dv=dx$, logo
  $du=\frac{1}{x}\,dx$, $v=x$.
  \[
    \int_1^e \ln x\,dx = \Big[x\ln x\Big]_1^e - \int_1^e x\cdot\frac{1}{x}\,dx
    = \Big[x\ln x\Big]_1^e - \Big[x\Big]_1^e
    = (e\cdot 1 - 1\cdot 0)-(e-1) = e - e + 1 = 1.
  \]

  \textbf{(b)} Substituição: $u=1-x^2$, $du=-2x\,dx$.
  Limites: $x=0\Rightarrow u=1$; $x=1\Rightarrow u=0$.
  \[
    \int_0^1 x\sqrt{1-x^2}\,dx = -\frac{1}{2}\int_1^0 \sqrt{u}\,du
    = \frac{1}{2}\int_0^1 u^{1/2}\,du
    = \frac{1}{2}\left[\frac{2}{3}u^{3/2}\right]_0^1 = \frac{1}{3}.
  \]

  \textbf{(c)} Identidade: $\sin^2 x = \dfrac{1-\cos 2x}{2}$.
  \[
    \int_0^{\pi/2}\sin^2 x\,dx = \int_0^{\pi/2}\frac{1-\cos 2x}{2}\,dx
    = \left[\frac{x}{2}-\frac{\sin 2x}{4}\right]_0^{\pi/2}
    = \frac{\pi}{4} - 0 = \frac{\pi}{4}.
  \]
\end{exercise}

\begin{exercise}
  Determine a área da região delimitada por $y=x^3$ e $y=x$.

  \medskip\textbf{Resolução.}

  \textbf{Passo 1 -- Pontos de interseção.} $x^3=x \Rightarrow x(x^2-1)=0
  \Rightarrow x\in\{-1,0,1\}$.

  \textbf{Passo 2 -- Sinal da diferença.} Em $(-1,0)$: $x^3-x = x(x^2-1)$.
  Teste em $x=-\tfrac{1}{2}$: $(-\tfrac{1}{2})({\tfrac{1}{4}-1})=(-\tfrac{1}{2})(-\tfrac{3}{4})>0$,
  logo $x^3\geq x$ em $(-1,0)$.
  Em $(0,1)$: teste em $x=\tfrac{1}{2}$: resultado negativo, logo $x\geq x^3$.

  \textbf{Passo 3 -- Cálculo por simetria.} A função $x^3-x$ é ímpar, pelo que
  a área em $(-1,0)$ iguala a área em $(0,1)$:
  \[
    A = 2\int_0^1(x-x^3)\,dx = 2\left[\frac{x^2}{2}-\frac{x^4}{4}\right]_0^1
    = 2\left(\frac{1}{2}-\frac{1}{4}\right) = 2\cdot\frac{1}{4} = \frac{1}{2}.
  \]
\end{exercise}

\begin{exercise}
  Determine o volume do sólido obtido ao rodar a região delimitada por
  $y=\sqrt{x}$, $x=4$ e $y=0$ em torno do eixo $Ox$.

  \medskip\textbf{Resolução.}

  \textbf{Passo 1 -- Identificar o método.} A região é delimitada por baixo por
  $y=0$ e por cima por $y=\sqrt{x}$, com $x\in[0,4]$. Aplica-se o método dos
  discos.

  \textbf{Passo 2 -- Aplicar a fórmula.}
  \[
    V = \pi\int_0^4 [\sqrt{x}]^2\,dx = \pi\int_0^4 x\,dx.
  \]

  \textbf{Passo 3 -- Calcular.}
  \[
    V = \pi\left[\frac{x^2}{2}\right]_0^4 = \pi\cdot\frac{16}{2} = 8\pi.
  \]
\end{exercise}

\begin{exercise}
  Calcule o comprimento de arco da curva $y = \ln(\cos x)$ de $x=0$ a
  $x=\pi/4$.

  \medskip\textbf{Resolução.}

  \textbf{Passo 1 -- Derivada.}
  \[
    y' = \frac{-\sin x}{\cos x} = -\tan x.
  \]

  \textbf{Passo 2 -- Expressão sob a raiz.}
  \[
    1+[y']^2 = 1+\tan^2 x = \sec^2 x.
  \]

  \textbf{Passo 3 -- Integral do comprimento.}
  \[
    L = \int_0^{\pi/4}\sqrt{\sec^2 x}\,dx = \int_0^{\pi/4}|\sec x|\,dx
    = \int_0^{\pi/4}\sec x\,dx.
  \]
  (Note-se que $\sec x > 0$ em $[0,\pi/4]$.)

  \textbf{Passo 4 -- Calcular $\int\sec x\,dx$.}
  Usando a primitiva conhecida $\int\sec x\,dx = \ln|\sec x+\tan x|+C$:
  \[
    L = \Big[\ln|\sec x + \tan x|\Big]_0^{\pi/4}
    = \ln\!\left(\sqrt{2}+1\right) - \ln(1+0)
    = \ln(\sqrt{2}+1).
  \]
\end{exercise}

% -----------------------------------
\subsection{Questões Sugeridas para a Prova Oral -- Capítulo 2}
% -----------------------------------

\begin{enumerate}
  \item Enuncie e demonstre o Teorema Fundamental do Cálculo Integral
        (Parte 2). Qual o papel da continuidade?
  \item Explique a diferença entre integral indefinido e integral definido.
        Por que razão a constante arbitrária $C$ desaparece no caso definido?
  \item Quando é que a primitivação por partes é mais útil? Dê um exemplo
        em que seja necessário aplicá-la duas vezes.
  \item Descreva geometricamente o que representa o integral definido
        $\int_a^b f(x)\,dx$ quando $f$ toma valores negativos em parte de $[a,b]$.
  \item Deduza a fórmula do comprimento de arco a partir de primeiros
        princípios, recorrendo ao Teorema de Pitágoras.
  \item Compare o método dos discos e o método das cascas para o cálculo
        de volumes. Em que circunstâncias se prefere cada um?
  \item Qual a relação entre a área sob uma curva e uma soma de Riemann?
        Descreva a convergência.
\end{enumerate}

\newpage
% ================================================================
%  CAPÍTULO 3 -- FUNÇÕES REAIS DE VÁRIAS VARIÁVEIS REAIS
% ================================================================
\section{Funções Reais de Várias Variáveis Reais}

Este capítulo abrange os tópicos do programa 3.1--3.5, incluindo domínio e
representação gráfica, limites e continuidade, derivadas parciais,
diferenciabilidade, regra da cadeia e extremos livres.

% -----------------------------------
\subsection{Definição, Domínio, Curvas de Nível e Representação Gráfica (3.1)}
% -----------------------------------

\begin{definition}[Função de Várias Variáveis]
  Uma \emph{função real de $n$ variáveis reais} é uma aplicação
  $f: D \subseteq \mathbb{R}^n \to \mathbb{R}$.
  Para $n=2$ escrevemos $f(x,y)$ e o seu \emph{gráfico} é a superfície
  \[
    \{(x,y,z)\in\mathbb{R}^3 \mid z = f(x,y),\, (x,y)\in D\}.
  \]
\end{definition}
 
\begin{definition}[Curva de Nível]
  Para $f:D\subseteq\mathbb{R}^2\to\mathbb{R}$ e uma constante $c\in\mathbb{R}$,
  a \emph{curva de nível} (linha de contorno) à altura $c$ é
  \[
    L_c = \{(x,y)\in D \mid f(x,y) = c\}.
  \]
\end{definition}
 
\begin{example}
  Descreva as curvas de nível de $f(x,y)=x^2+y^2$.
 
  \textbf{Resolução.} A equação $f(x,y)=c$ escreve-se $x^2+y^2=c$.
  \begin{itemize}[nosep]
    \item \textbf{Passo 1:} Se $c<0$, a equação $x^2+y^2=c$ não tem solução
          real (soma de quadrados nunca é negativa), pelo que a curva de nível
          é vazia.
    \item \textbf{Passo 2:} Se $c=0$, a única solução é $x=y=0$, ou seja,
          a curva de nível é o ponto $(0,0)$.
    \item \textbf{Passo 3:} Se $c>0$, a equação $x^2+y^2=c$ representa
          uma circunferência centrada na origem com raio $\sqrt{c}$.
  \end{itemize}
  Conclusão: as curvas de nível são circunferências concêntricas de raio
  $\sqrt{c}$ (para $c>0$), o ponto origem ($c=0$), ou vazias ($c<0$).
\end{example}
 
\subsubsection*{Representação Gráfica}
 
O gráfico de $f(x,y)$ é uma \emph{superfície} no espaço $\mathbb{R}^3$.
Para o representar, os métodos mais comuns são:
 
\begin{itemize}[nosep]
  \item \textbf{Curvas de nível:} traçam-se as curvas $f(x,y)=c$ no plano
        $Oxy$ para vários valores de $c$, obtendo uma vista "de cima" da
        superfície. Curvas muito próximas indicam grande variação de $f$;
        curvas afastadas indicam variação suave.
  \item \textbf{Cortes verticais:} fixa-se $x=a$ ou $y=b$ e traça-se a
        curva resultante em $\mathbb{R}^2$, obtendo uma "fatia" da superfície.
\end{itemize}
 
\begin{example}[Representação Gráfica]
  Descreva o gráfico de $f(x,y)=4-x^2-y^2$.
 
  \textbf{Resolução.}
 
  \textbf{Passo 1 -- Identificar o tipo de superfície.} A equação
  $z=4-x^2-y^2$ pode escrever-se $x^2+y^2+z=4$, ou seja,
  $x^2+y^2=4-z$. Trata-se de um \emph{parabolóide de revolução} com
  vértice em $(0,0,4)$, aberto para baixo.
 
  \textbf{Passo 2 -- Curvas de nível.} Para $z=c$:
  \[
    x^2+y^2 = 4-c.
  \]
  São circunferências de raio $\sqrt{4-c}$, definidas para $c\leq 4$.
  Quando $c=4$ obtém-se o ponto $(0,0)$; quando $c\to-\infty$ o raio
  cresce indefinidamente.
 
  \textbf{Passo 3 -- Cortes verticais.} Fixando $y=0$:
  $z=4-x^2$, que é uma parábola no plano $Oxz$. Por simetria de
  revolução, qualquer corte vertical pelo eixo $Oz$ dá a mesma parábola.
\end{example}
 
\begin{example}[Domínio]
  Determine o domínio de $f(x,y)=\dfrac{\ln(1-x^2-y^2)}{\sqrt{x+y}}$.
 
  \textbf{Resolução.} É necessário:
  \begin{itemize}[nosep]
    \item $1-x^2-y^2>0 \Rightarrow x^2+y^2 < 1$ (disco unitário aberto).
    \item $x+y>0$ (acima da reta $y=-x$).
  \end{itemize}
  Assim $D = \{(x,y)\mid x^2+y^2<1 \text{ e } x+y>0\}$.
\end{example}

% -----------------------------------
\subsection{Limites e Continuidade (3.2)}
% -----------------------------------

\begin{definition}[Limite]
  Diz-se que $\displaystyle\lim_{(x,y)\to(a,b)}f(x,y) = L$ se para todo o
  $\varepsilon>0$ existe $\delta>0$ tal que
  \[
    0 < \sqrt{(x-a)^2+(y-b)^2} < \delta \implies |f(x,y)-L|<\varepsilon.
  \]
\end{definition}

\begin{remark}
  Ao contrário de $\mathbb{R}$, em $\mathbb{R}^2$ o limite tem de ser o mesmo
  ao longo de \emph{todos} os caminhos. Uma estratégia comum para mostrar que
  um limite \emph{não existe} consiste em exibir dois caminhos que dão valores
  diferentes.
\end{remark}

\begin{example}[Limite inexistente]
  Mostre que $\displaystyle\lim_{(x,y)\to(0,0)}\frac{xy}{x^2+y^2}$ não existe.

  \textbf{Resolução.}
  Pelo caminho $y=0$: $f(x,0)=0\to 0$.
  Pelo caminho $y=x$: $f(x,x)=\frac{x^2}{2x^2}=\frac{1}{2}\to\frac{1}{2}$.
  Como os limites são diferentes, o limite não existe.
\end{example}

\begin{definition}[Continuidade]
  $f$ é \emph{contínua} em $(a,b)$ se
  $\displaystyle\lim_{(x,y)\to(a,b)}f(x,y)=f(a,b)$.
\end{definition}

% -----------------------------------
\subsection{Derivadas Parciais e Diferenciabilidade (3.3)}
% -----------------------------------

\begin{definition}[Derivada Parcial]
  A \emph{derivada parcial} de $f$ em ordem a $x$ no ponto $(a,b)$ é
  \[
    \frac{\partial f}{\partial x}(a,b)
    = \lim_{h\to 0}\frac{f(a+h,b)-f(a,b)}{h},
  \]
  calculada derivando $f$ em ordem a $x$ mantendo $y$ constante.
\end{definition}

\begin{definition}[Gradiente]
  O \emph{gradiente} de $f$ em $(a,b)$ é o vetor
  \[
    \nabla f(a,b) = \left(\frac{\partial f}{\partial x}(a,b),\,
                          \frac{\partial f}{\partial y}(a,b)\right).
  \]
  Aponta na direção de maior crescimento de $f$.
\end{definition}

\begin{definition}[Derivada Direcional]
  Para um vetor unitário $\mathbf{u}=(u_1,u_2)$, a derivada direcional é
  \[
    D_{\mathbf{u}}f(a,b) = \nabla f(a,b)\cdot\mathbf{u}
      = \frac{\partial f}{\partial x}u_1 + \frac{\partial f}{\partial y}u_2.
  \]
\end{definition}

\begin{definition}[Diferenciabilidade]
  $f$ é \emph{diferenciável} em $(a,b)$ se existir uma aplicação linear
  $L:\mathbb{R}^2\to\mathbb{R}$ tal que
  \[
    \lim_{(h,k)\to(0,0)}\frac{f(a+h,b+k)-f(a,b)-L(h,k)}{\sqrt{h^2+k^2}}=0.
  \]
  A aplicação linear é $L(h,k)=\nabla f(a,b)\cdot(h,k)$.
\end{definition}

\begin{keybox}
\textbf{Condição suficiente de diferenciabilidade:} Se $f_x$ e $f_y$
existirem e forem contínuas numa vizinhança de $(a,b)$, então $f$ é
diferenciável em $(a,b)$.

\medskip
\textbf{Cadeia de implicações:}
Diferenciável $\Rightarrow$ Contínua $\Rightarrow$ Limite existe;
mas Derivadas parciais existem $\not\Rightarrow$ Diferenciável.
\end{keybox}

\begin{example}
  Determine $\nabla f$ para $f(x,y)=x^3 y + \sin(xy)$.

  \textbf{Resolução.}
  \[
    f_x = 3x^2 y + y\cos(xy), \qquad f_y = x^3 + x\cos(xy).
  \]
\end{example}

\begin{example}[Derivadas parciais de segunda ordem -- Hessiana]
  Para $f(x,y)=e^{x^2+y^2}$, calcule todas as derivadas parciais de
  segunda ordem.

  \textbf{Resolução.} $f_x = 2xe^{x^2+y^2}$, $f_y = 2ye^{x^2+y^2}$.
  \[
    f_{xx}=(2+4x^2)e^{x^2+y^2},\quad f_{yy}=(2+4y^2)e^{x^2+y^2},\quad
    f_{xy}=f_{yx}=4xye^{x^2+y^2}.
  \]
  (Note-se que $f_{xy}=f_{yx}$ pelo Teorema de Schwarz/Clairaut, uma vez
  que todas as derivadas parciais de segunda ordem são contínuas.)
\end{example}

% -----------------------------------
\subsection{Regra da Cadeia (3.4)}
% -----------------------------------

\begin{theorem}[Regra da Cadeia]
  Sejam $z=f(x,y)$, $x=x(t)$, $y=y(t)$ diferenciáveis. Então
  \[
    \frac{dz}{dt} = \frac{\partial f}{\partial x}\frac{dx}{dt}
                  + \frac{\partial f}{\partial y}\frac{dy}{dt}.
  \]
  De modo mais geral, se $x=x(s,t)$, $y=y(s,t)$:
  \[
    \frac{\partial z}{\partial s}
      = \frac{\partial f}{\partial x}\frac{\partial x}{\partial s}
      + \frac{\partial f}{\partial y}\frac{\partial y}{\partial s}.
  \]
\end{theorem}

\begin{example}
  Seja $z = x^2+y^2$, $x=\cos t$, $y=\sin t$. Calcule $\frac{dz}{dt}$.

  \textbf{Resolução.}
  \[
    \frac{dz}{dt} = 2x(-\sin t)+2y(\cos t) = -2\cos t\sin t + 2\sin t\cos t = 0.
  \]
  (Faz sentido: $z = \cos^2 t + \sin^2 t = 1$.)
\end{example}

% -----------------------------------
\subsection{Extremos Livres (3.5)}
% -----------------------------------

\begin{definition}[Ponto Crítico]
  $(a,b)$ é um \emph{ponto crítico} de $f$ se $\nabla f(a,b)=\mathbf{0}$,
  ou seja, $f_x(a,b)=0$ e $f_y(a,b)=0$.
\end{definition}

\begin{theorem}[Teste da Segunda Derivada]
  Seja $(a,b)$ um ponto crítico de $f$. Define-se o \emph{determinante
  Hessiano}
  \[
    H = f_{xx}(a,b)\,f_{yy}(a,b) - [f_{xy}(a,b)]^2.
  \]
  \begin{itemize}[nosep]
    \item $H > 0$ e $f_{xx}(a,b)>0$: \textbf{mínimo local}.
    \item $H > 0$ e $f_{xx}(a,b)<0$: \textbf{máximo local}.
    \item $H < 0$: \textbf{ponto de sela}.
    \item $H = 0$: o teste é \textbf{inconclusivo}.
  \end{itemize}
\end{theorem}

\begin{example}
  Determine e classifique os pontos críticos de
  $f(x,y)=x^3-3x+y^2-2y$.

  \textbf{Resolução.}
  \[
    f_x = 3x^2-3 = 0 \implies x=\pm 1,\qquad
    f_y = 2y-2 = 0 \implies y = 1.
  \]
  Pontos críticos: $(1,1)$ e $(-1,1)$.

  Derivadas de segunda ordem: $f_{xx}=6x$, $f_{yy}=2$, $f_{xy}=0$.

  \begin{center}
  \begin{tabular}{cccccc}
    \toprule
    Ponto & $f_{xx}$ & $f_{yy}$ & $f_{xy}$ & $H$ & Classificação \\
    \midrule
    $(1,1)$  & $6$  & $2$ & $0$ & $12>0$, $f_{xx}>0$ & Mínimo local \\
    $(-1,1)$ & $-6$ & $2$ & $0$ & $-12<0$ & Ponto de sela \\
    \bottomrule
  \end{tabular}
  \end{center}

  Valor mínimo local: $f(1,1) = 1-3+1-2 = -3$.
\end{example}

% -----------------------------------
\subsection{Exercícios de Prática -- Capítulo 3}
% -----------------------------------

\begin{exercise}
  Para $f(x,y)=\sqrt{9-x^2-y^2}$, determine o domínio e descreva as
  curvas de nível.

  \medskip\textbf{Resolução.}

  \textbf{Passo 1 -- Domínio.} A raiz quadrada exige que o radicando seja
  não-negativo:
  \[
    9-x^2-y^2 \geq 0 \iff x^2+y^2 \leq 9.
  \]
  O domínio é o disco fechado $D = \{(x,y)\in\mathbb{R}^2 \mid x^2+y^2\leq 9\}$,
  de raio $3$ centrado na origem.

  \textbf{Passo 2 -- Curvas de nível.} Para $c\geq 0$, a equação $f(x,y)=c$
  equivale a $\sqrt{9-x^2-y^2}=c$, ou seja,
  \[
    x^2+y^2 = 9-c^2, \quad \text{com } 0\leq c\leq 3.
  \]
  \begin{itemize}[nosep]
    \item Se $0\leq c < 3$: circunferência de raio $\sqrt{9-c^2}$.
    \item Se $c=3$: o ponto $(0,0)$ (raio nulo).
    \item Se $c>3$ ou $c<0$: curva vazia (fora do contradomínio).
  \end{itemize}
  As curvas de nível são portanto circunferências concêntricas de raio
  decrescente à medida que $c$ aumenta de $0$ a $3$, correspondendo
  geometricamente a paralelos de uma semiesfera de raio $3$.
\end{exercise}

\begin{exercise}
  Averigue se os seguintes limites existem e calcule-os caso existam:
  \begin{enumerate}[label=(\alph*)]
    \item $\displaystyle\lim_{(x,y)\to(0,0)}\frac{x^2 y}{x^4+y^2}$
    \item $\displaystyle\lim_{(x,y)\to(1,2)}(x^2+3xy)$
    \item $\displaystyle\lim_{(x,y)\to(0,0)}\frac{x^3+y^3}{x^2+y^2}$
  \end{enumerate}

  \medskip\textbf{Resolução.}

  \textbf{(a)} Testa-se com dois caminhos distintos.

  \emph{Caminho 1:} $y=0$. $f(x,0)=\dfrac{0}{x^4}=0\to 0$.

  \emph{Caminho 2:} $y=x^2$. $f(x,x^2)=\dfrac{x^2\cdot x^2}{x^4+x^4}=\dfrac{x^4}{2x^4}=\dfrac{1}{2}\to\dfrac{1}{2}$.

  Como os dois limites são diferentes ($0\neq\tfrac{1}{2}$), o limite
  \textbf{não existe}.

  \textbf{(b)} A função $g(x,y)=x^2+3xy$ é polinomial, portanto contínua em
  todo o $\mathbb{R}^2$. Assim:
  \[
    \lim_{(x,y)\to(1,2)}(x^2+3xy) = 1^2+3\cdot 1\cdot 2 = 1+6 = 7.
  \]

  \textbf{(c)} Passa-se para coordenadas polares: $x=r\cos\theta$,
  $y=r\sin\theta$, com $r\to 0^+$.
  \[
    \frac{x^3+y^3}{x^2+y^2}
    = \frac{r^3(\cos^3\theta+\sin^3\theta)}{r^2}
    = r(\cos^3\theta+\sin^3\theta).
  \]
  Como $|\cos^3\theta+\sin^3\theta|\leq 2$ para todo o $\theta$, temos
  \[
    \left|r(\cos^3\theta+\sin^3\theta)\right|\leq 2r\to 0.
  \]
  Portanto o limite \textbf{existe e é igual a $0$}.
\end{exercise}

\begin{exercise}
  Calcule todas as derivadas parciais de primeira e segunda ordem de
  $f(x,y) = x^2\sin y + y\ln x$.

  \medskip\textbf{Resolução.}

  \textbf{Passo 1 -- Derivadas de primeira ordem.}
  \begin{align*}
    f_x &= \frac{\partial}{\partial x}(x^2\sin y) + \frac{\partial}{\partial x}(y\ln x)
         = 2x\sin y + \frac{y}{x}.\\[4pt]
    f_y &= \frac{\partial}{\partial y}(x^2\sin y) + \frac{\partial}{\partial y}(y\ln x)
         = x^2\cos y + \ln x.
  \end{align*}

  \textbf{Passo 2 -- Derivadas de segunda ordem.}
  \begin{align*}
    f_{xx} &= \frac{\partial}{\partial x}\!\left(2x\sin y+\frac{y}{x}\right)
            = 2\sin y - \frac{y}{x^2}.\\[4pt]
    f_{yy} &= \frac{\partial}{\partial y}(x^2\cos y+\ln x)
            = -x^2\sin y.\\[4pt]
    f_{xy} &= \frac{\partial}{\partial y}\!\left(2x\sin y+\frac{y}{x}\right)
            = 2x\cos y + \frac{1}{x}.\\[4pt]
    f_{yx} &= \frac{\partial}{\partial x}(x^2\cos y+\ln x)
            = 2x\cos y + \frac{1}{x}.
  \end{align*}
  Confirma-se que $f_{xy}=f_{yx}$ (Teorema de Schwarz), pois ambas as
  derivadas mistas são contínuas no domínio $x>0$.
\end{exercise}

\begin{exercise}
  Seja $f(x,y) = x^2 y - y^3$. Determine a derivada direcional em $(1,1)$
  na direção de $\mathbf{v}=(3,4)$.

  \medskip\textbf{Resolução.}

  \textbf{Passo 1 -- Normalizar o vetor.}
  \[
    \|\mathbf{v}\| = \sqrt{3^2+4^2} = \sqrt{25} = 5,\qquad
    \mathbf{u} = \frac{\mathbf{v}}{\|\mathbf{v}\|} = \left(\frac{3}{5},\frac{4}{5}\right).
  \]

  \textbf{Passo 2 -- Calcular o gradiente.}
  \[
    f_x = 2xy,\quad f_y = x^2 - 3y^2.
  \]
  Em $(1,1)$: $f_x(1,1)=2$, $f_y(1,1)=1-3=-2$.
  Logo $\nabla f(1,1)=(2,-2)$.

  \textbf{Passo 3 -- Produto interno.}
  \[
    D_{\mathbf{u}}f(1,1) = \nabla f(1,1)\cdot\mathbf{u}
    = 2\cdot\frac{3}{5}+(-2)\cdot\frac{4}{5}
    = \frac{6}{5}-\frac{8}{5} = -\frac{2}{5}.
  \]
  A derivada direcional é $-\dfrac{2}{5}$, ou seja, $f$ decresce nessa
  direção.
\end{exercise}

\begin{exercise}
  Dado $z = e^{u^2+v^2}$, $u = x\cos y$, $v = x\sin y$, utilize a regra
  da cadeia para calcular $\partial z/\partial x$ e $\partial z/\partial y$.

  \medskip\textbf{Resolução.}

  \textbf{Passo 1 -- Derivadas intermédias.}
  \[
    \frac{\partial z}{\partial u}=2ue^{u^2+v^2},\quad
    \frac{\partial z}{\partial v}=2ve^{u^2+v^2}.
  \]
  \[
    \frac{\partial u}{\partial x}=\cos y,\quad
    \frac{\partial u}{\partial y}=-x\sin y,\quad
    \frac{\partial v}{\partial x}=\sin y,\quad
    \frac{\partial v}{\partial y}=x\cos y.
  \]

  \textbf{Passo 2 -- Regra da cadeia para $\partial z/\partial x$.}
  \[
    \frac{\partial z}{\partial x}
    = \frac{\partial z}{\partial u}\frac{\partial u}{\partial x}
    + \frac{\partial z}{\partial v}\frac{\partial v}{\partial x}
    = 2ue^{u^2+v^2}\cos y + 2ve^{u^2+v^2}\sin y
    = 2e^{u^2+v^2}(u\cos y + v\sin y).
  \]

  \textbf{Passo 3 -- Simplificar.} Substituindo $u=x\cos y$ e $v=x\sin y$:
  \[
    u\cos y+v\sin y = x\cos^2 y + x\sin^2 y = x.
  \]
  Como $u^2+v^2=x^2\cos^2 y+x^2\sin^2 y=x^2$:
  \[
    \frac{\partial z}{\partial x} = 2xe^{x^2}.
  \]

  \textbf{Passo 4 -- Regra da cadeia para $\partial z/\partial y$.}
  \[
    \frac{\partial z}{\partial y}
    = 2ue^{u^2+v^2}(-x\sin y) + 2ve^{u^2+v^2}(x\cos y)
    = 2xe^{x^2}(-u\sin y+v\cos y).
  \]
  Substituindo: $-u\sin y+v\cos y = -x\cos y\sin y+x\sin y\cos y = 0$.
  \[
    \frac{\partial z}{\partial y} = 0.
  \]
  (Faz sentido: $z=e^{x^2}$ não depende de $y$.)
\end{exercise}

\begin{exercise}
  Determine e classifique todos os pontos críticos de:
  \begin{enumerate}[label=(\alph*)]
    \item $f(x,y) = x^2 + xy + y^2 - 3x$
    \item $f(x,y) = x^3 + y^3 - 3xy$
    \item $f(x,y) = \sin x + \sin y + \sin(x+y)$, em
          $(0,\pi/2)\times(0,\pi/2)$
  \end{enumerate}

  \medskip\textbf{Resolução.}

  \textbf{(a)} $f(x,y) = x^2+xy+y^2-3x$.

  \emph{Passo 1 -- Condições de primeira ordem.}
  \[
    f_x = 2x+y-3=0,\qquad f_y = x+2y=0.
  \]
  De $f_y=0$: $x=-2y$. Substituindo em $f_x=0$: $-4y+y-3=0\Rightarrow y=-1$,
  $x=2$.

  Ponto crítico: $(2,-1)$.

  \emph{Passo 2 -- Teste da segunda derivada.}
  $f_{xx}=2$, $f_{yy}=2$, $f_{xy}=1$.
  \[
    H = 2\cdot 2 - 1^2 = 3 > 0,\quad f_{xx}=2>0.
  \]
  Conclusão: $(2,-1)$ é um \textbf{mínimo local}, com $f(2,-1)=4-2+1-6=-3$.

  \bigskip
  \textbf{(b)} $f(x,y)=x^3+y^3-3xy$.

  \emph{Passo 1 -- Condições de primeira ordem.}
  \[
    f_x = 3x^2-3y=0 \Rightarrow y=x^2,\qquad
    f_y = 3y^2-3x=0 \Rightarrow x=y^2.
  \]
  Substituindo $y=x^2$ em $x=y^2$: $x=x^4\Rightarrow x(x^3-1)=0$, logo
  $x=0$ ou $x=1$.
  \begin{itemize}[nosep]
    \item $x=0$: $y=0$. Ponto $(0,0)$.
    \item $x=1$: $y=1$. Ponto $(1,1)$.
  \end{itemize}

  \emph{Passo 2 -- Teste da segunda derivada.}
  $f_{xx}=6x$, $f_{yy}=6y$, $f_{xy}=-3$.

  \begin{center}
  \begin{tabular}{cccccc}
    \toprule
    Ponto & $f_{xx}$ & $f_{yy}$ & $f_{xy}$ & $H$ & Classificação \\
    \midrule
    $(0,0)$ & $0$ & $0$ & $-3$ & $0-9=-9<0$ & Ponto de sela \\
    $(1,1)$ & $6$ & $6$ & $-3$ & $36-9=27>0$, $f_{xx}>0$ & Mínimo local \\
    \bottomrule
  \end{tabular}
  \end{center}

  Valor mínimo local: $f(1,1)=1+1-3=-1$.

  \bigskip
  \textbf{(c)} $f(x,y)=\sin x+\sin y+\sin(x+y)$, em $(0,\tfrac{\pi}{2})\times(0,\tfrac{\pi}{2})$.

  \emph{Passo 1 -- Condições de primeira ordem.}
  \[
    f_x = \cos x + \cos(x+y) = 0,\qquad
    f_y = \cos y + \cos(x+y) = 0.
  \]
  Subtraindo: $\cos x - \cos y = 0 \Rightarrow \cos x = \cos y$.
  No domínio $(0,\tfrac{\pi}{2})$ o cosseno é estritamente decrescente, logo
  $x=y$.

  Substituindo $y=x$ em $f_x=0$:
  $\cos x + \cos 2x = 0$.
  Usando $\cos 2x = 2\cos^2 x - 1$:
  \[
    2\cos^2 x + \cos x - 1 = 0 \Rightarrow (2\cos x - 1)(\cos x + 1) = 0.
  \]
  No domínio aberto $(0,\tfrac{\pi}{2})$: $\cos x = \tfrac{1}{2}
  \Rightarrow x=\tfrac{\pi}{3}$.

  Ponto crítico: $\left(\tfrac{\pi}{3},\tfrac{\pi}{3}\right)$.

  \emph{Passo 2 -- Teste da segunda derivada.}
  \[
    f_{xx}=-\sin x-\sin(x+y),\quad
    f_{yy}=-\sin y-\sin(x+y),\quad
    f_{xy}=-\sin(x+y).
  \]
  Em $\left(\tfrac{\pi}{3},\tfrac{\pi}{3}\right)$: $x+y=\tfrac{2\pi}{3}$,
  $\sin\tfrac{\pi}{3}=\tfrac{\sqrt{3}}{2}$, $\sin\tfrac{2\pi}{3}=\tfrac{\sqrt{3}}{2}$.
  \[
    f_{xx}=f_{yy}=-\frac{\sqrt{3}}{2}-\frac{\sqrt{3}}{2}=-\sqrt{3},\quad
    f_{xy}=-\frac{\sqrt{3}}{2}.
  \]
  \[
    H = (-\sqrt{3})^2-\left(-\frac{\sqrt{3}}{2}\right)^2 = 3-\frac{3}{4}=\frac{9}{4}>0,
    \quad f_{xx}=-\sqrt{3}<0.
  \]
  Conclusão: $\left(\tfrac{\pi}{3},\tfrac{\pi}{3}\right)$ é um \textbf{máximo local},
  com $f=\sin\tfrac{\pi}{3}+\sin\tfrac{\pi}{3}+\sin\tfrac{2\pi}{3}=\tfrac{3\sqrt{3}}{2}$.
\end{exercise}

% -----------------------------------
\subsection{Questões Sugeridas para a Prova Oral -- Capítulo 3}
% -----------------------------------

\begin{enumerate}
  \item Explique por que razão a existência de derivadas parciais num ponto
        \emph{não} garante a continuidade nesse ponto. Dê um contra-exemplo.
  \item Enuncie as condições em que as derivadas parciais mistas $f_{xy}$
        e $f_{yx}$ são iguais (Teorema de Schwarz/Clairaut).
  \item Interprete geometricamente o gradiente: o que nos diz $\nabla f(a,b)$
        sobre a superfície $z=f(x,y)$?
  \item Deduza a fórmula da regra da cadeia para $dz/dt$ quando
        $z=f(x(t),y(t))$ a partir da definição de diferenciabilidade.
  \item Descreva o teste da segunda derivada. O que acontece quando o
        determinante Hessiano é zero? Apresente exemplos de cada caso.
  \item Como generaliza o conceito de derivada direcional a noção de
        derivada parcial? Quando é que $D_{\mathbf{u}}f$ é maximizada, e em
        que direção?
  \item Explique a ligação entre curvas de nível, o gradiente e a direção
        de maior subida/descida.
\end{enumerate}

\end{document}